\documentclass[11pt]{article}
\usepackage[T1]{fontenc}
\usepackage[utf8]{inputenc}
\usepackage[spanish]{babel}
\usepackage{amsmath,amssymb}
\usepackage[a4paper,margin=2.3cm]{geometry}
\usepackage{graphicx}
\usepackage{tikz}
\usepackage{float}
\usetikzlibrary{arrows.meta,babel}
\tikzset{
  repaso/.style={>=Stealth,line cap=round,line join=round,
                 every node/.style={font=\small}},
  borde/.style={black,line width=.72pt},
  eje/.style={-{Stealth[length=1.45mm]},thin},
  vector/.style={-{Stealth[length=1.85mm]},line width=.75pt},
  trayecto/.style={-{Stealth[length=1.6mm]},line width=.65pt},
  guia/.style={densely dashed,black!66,line width=.5pt},
  medida/.style={{Stealth[length=1.45mm]}-{Stealth[length=1.45mm]},thin},
  llamada/.style={black!65,line width=.38pt},
  corte/.style={fill=white,inner sep=.55pt}
}
\spanishdecimal{.}
\setlength{\parindent}{0pt}
\setlength{\parskip}{0.35em}
\setlength{\abovedisplayskip}{4pt}
\setlength{\belowdisplayskip}{4pt}
\setlength{\abovedisplayshortskip}{4pt}
\setlength{\belowdisplayshortskip}{4pt}

\newcommand{\encabezado}[3]{%
  {\Large\bfseries #1\par}
  \vspace{2pt}
  {\normalsize #2\hfill #3\par}
  \vspace{6pt}\hrule\vspace{8pt}}

\newcommand{\problema}[2]{%
  \vspace{0.55em}
  \noindent\textbf{#1.}\enspace{\small #2}\par
  \vspace{0.3em}}

\begin{document}



\encabezado{Repaso E2 -- Problemas resueltos}
{Electrodinámica I}{P.S., J.A., J.R.}

\section*{Enunciados de los Problemas}


\section*{Problema 1: Electrostática en la Materia (Esfera Conductora con Corteza Dieléctrica)}

Una esfera metálica conductora maciza de radio $a$ posee una carga libre neta $+Q$. La esfera está rodeada concéntricamente por un cascarón dieléctrico lineal, homogéneo e isotrópico de permitividad eléctrica $\varepsilon$, el cual ocupa la región comprendida entre el radio interno $a$ y el radio externo $b$ ($b > a$). Para $r > b$ el sistema se encuentra rodeado de vacío (permitividad $\varepsilon_0$).

\begin{enumerate}
    \item[a)] Utilizando la simetría esférica y la Ley de Gauss para medios materiales, determine los vectores desplazamiento eléctrico $\mathbf{D}(r)$, campo eléctrico $\mathbf{E}(r)$ y polarización $\mathbf{P}(r)$ en las tres regiones del espacio: $r < a$, $a < r < b$, y $r > b$.
    \item[b)] Tomando como referencia de potencial nulo el infinito ($V(\infty) = 0$), calcule el potencial electrostático $V(a)$ de la esfera conductora y deduzca la capacitancia $C$ de este conductor recubierto. Evalúe e interprete físicamente los casos límite $\varepsilon \to \varepsilon_0$ y $b \to \infty$.
    \item[c)] Determine la densidad volumétrica de carga ligada $\rho_b$ en el interior del dieléctrico ($a < r < b$) y las densidades superficiales de carga ligada $\sigma_b$ en las fronteras $r = a$ y $r = b$. Compruebe explícitamente que la carga ligada total del cascarón es cero y que la suma de la carga libre $+Q$ con la carga ligada en $r = a$ explica el apantallamiento parcial del campo eléctrico dentro del dieléctrico.
\end{enumerate}

\vspace{0.3cm}

\section*{Problema 2: Magnetostática en el Vacío (Cilindro Cargado en Rotación)}

Un cilindro hueco, muy largo y de paredes aislantes muy delgadas, tiene radio $R$ y posee una densidad superficial de carga eléctrica uniforme $\sigma$ adherida rígidamente a su manto. El cilindro gira alrededor de su eje de simetría (el eje $z$) con velocidad angular constante $\boldsymbol{\omega} = \omega \hat{\mathbf{z}}$.

\begin{enumerate}
    \item[a)] Determine el vector densidad de corriente superficial $\mathbf{K}$ que aparece sobre la pared del cilindro ($s = R$) producto de la rotación.
    \item[b)] Utilizando argumentos de simetría cilíndrica y la Ley de Ampère en forma integral (asumiendo que el campo magnético se anula a distancias radiales infinitas), calcule el vector campo magnético $\mathbf{B}(s)$ en el interior ($s < R$) y en el exterior ($s > R$) del cilindro. Verifique explícitamente que su resultado cumple la condición de frontera $\mathbf{B}_{\text{ext}} - \mathbf{B}_{\text{int}} = \mu_0 (\mathbf{K} \times \hat{\mathbf{n}})$.
    \item[c)] A partir del campo magnético obtenido, encuentre el potencial vectorial magnético continuo $\mathbf{A}(s)$ en todo el espacio imponiendo $\mathbf{A}(0) = \mathbf{0}$, compruebe que satisface el gauge de Coulomb ($\nabla \cdot \mathbf{A} = 0$), y calcule el momento dipolar magnético por unidad de longitud axial del cilindro.
\end{enumerate}

\vspace{0.3cm}

\section*{Problema 3: Magnetostática en la Materia e Inducción (Toroide con Entrehierro)}

Un núcleo toroidal de radio medio $R$ y sección transversal pequeña de área $A$ ($\sqrt{A} \ll R$, de modo que la magnitud de los campos puede aproximarse como uniforme sobre la sección transversal) está construido con un material magnético lineal de permeabilidad $\mu = \mu_0(1 + \chi_m)$. Al anillo toroidal se le practica un corte transversal muy estrecho de espesor $d$ ($d \ll R$), dejando una ranura de vacío (``entrehierro''). Alrededor del núcleo se enrolla un devanado primario de $N$ espiras por el cual circula una corriente libre $I$.

\begin{enumerate}
    \item[a)] Despreciando el ensanchamiento de las líneas de campo en los bordes de la ranura, aplique la continuidad de la componente normal de $\mathbf{B}$ en las caras del entrehierro y la Ley de Ampère para el campo auxiliar $\mathbf{H}$ a lo largo de la circunferencia media de longitud $2\pi R$. Determine la magnitud del campo magnético $B$ en el circuito magnético, así como las magnitudes $H_m$ en el material y $H_0$ en la ranura de vacío.
    \item[b)] Calcule la magnitud de la densidad superficial de corriente de magnetización $K_b$ que circula sobre el manto lateral del núcleo magnético. Explique físicamente por qué $H_0$ y $H_m$ difieren drásticamente al cruzar las caras del entrehierro pese a que no existen corrientes libres en dichas superficies.
    \item[c)] Si se enrolla una bobina secundaria de $N_2$ vueltas sobre el mismo núcleo y se alimenta el devanado primario con una corriente alterna $I(t) = I_0 \cos(\omega t)$ (en régimen cuasiestático), determine el coeficiente de inductancia mutua $\mathcal{M}$ entre ambas bobinas y la fuerza electromotriz $\mathcal{E}_2(t)$ inducida en el secundario según la Ley de Faraday.
\end{enumerate}

\vspace{0.3cm}

\section*{Problema 4: Ley de Ampère-Maxwell y Ley de Faraday en Medios Materiales}

Este problema contrasta dos configuraciones cilíndricas no estacionarias en presencia de materia:

\begin{enumerate}
    \item[a)] \textbf{Descarga óhmica radial:} Dos cilindros conductores coaxiales muy largos de radios $a$ y $b$ ($b > a$) tienen el espacio intermedio relleno de un medio material lineal, homogéneo e isotrópico con permitividad $\varepsilon$, permeabilidad $\mu_0$ y conductividad eléctrica débil $\sigma$ (donde rige $\mathbf{J}_f = \sigma \mathbf{E}$), sin cables externos conectados. En $t = 0$, el cilindro interno posee una densidad lineal de carga libre $+\lambda_0$ y el externo $-\lambda_0$. Utilizando la Ley de Gauss para $\mathbf{D}$ y la conservación de la carga libre, encuentre $\mathbf{D}(s,t)$, $\mathbf{E}(s,t)$ y $\mathbf{J}_f(s,t)$ para $a < s < b$ y $t > 0$, demostrando que $\lambda(t) = \lambda_0 e^{-t/\tau}$ con $\tau = \varepsilon / \sigma$.
    \item[b)] \textbf{Cancelación magnética exacta:} Para el sistema del inciso anterior, calcule la densidad de corriente de desplazamiento $\mathbf{J}_d = \frac{\partial \mathbf{D}}{\partial t}$ en $a < s < b$ y aplique la Ley de Ampère-Maxwell en la materia ($\oint \mathbf{H} \cdot d\mathbf{l} = I_{f,\text{enc}} + \frac{d\Phi_D}{dt}$) para demostrar que el campo magnético $\mathbf{B}$ inducido durante toda la descarga es idénticamente nulo en todo el espacio.
\end{enumerate}

\vspace{0.3cm}

\newpage

\section*{Problema 5: Ecuaciones de Maxwell en la Materia (Onda Estacionaria y Fronteras)}

En el interior de un medio dieléctrico y magnético lineal, homogéneo, isotrópico y sin pérdidas ($\rho_f = 0$, $\mathbf{J}_f = \mathbf{0}$, con permitividad $\varepsilon = \varepsilon_r \varepsilon_0$ y permeabilidad $\mu = \mu_r \mu_0$), existe un campo magnético oscilante descrito en coordenadas cartesianas por:
\begin{equation}
    \mathbf{B}(x, y, z, t) = B_0 \cos(kz) \sin(\omega t) \hat{\mathbf{y}}
\end{equation}
donde $B_0$, $k$ y $\omega$ son constantes positivas.

\begin{enumerate}
    \item[a)] Compruebe que $\mathbf{B}$ satisface $\nabla \cdot \mathbf{B} = 0$, y escriba explícitamente el campo auxiliar $\mathbf{H}(z,t)$ y el vector magnetización $\mathbf{M}(z,t)$ en el medio.
    \item[b)] Aplicando la ecuación diferencial de Ampère-Maxwell en la materia ($\nabla \times \mathbf{H} = \frac{\partial \mathbf{D}}{\partial t}$) y asumiendo ausencia de campos electrostáticos estacionarios, determine los vectores $\mathbf{D}(z,t)$, $\mathbf{E}(z,t)$ y la polarización $\mathbf{P}(z,t)$. Verifique que se cumple la Ley de Gauss macroscópica ($\nabla \cdot \mathbf{D} = 0$).
    \item[c)] Sustituya los campos en la Ley de Faraday ($\nabla \times \mathbf{E} = -\frac{\partial \mathbf{B}}{\partial t}$) para deducir la relación algebraica exacta que vincula $\omega$, $k$, $\varepsilon_r$, $\mu_r$ y la rapidez de la luz en el vacío $c = 1/\sqrt{\varepsilon_0 \mu_0}$. Finalmente, si este medio material está confinado entre dos placas conductoras perfectas situadas en $z = 0$ y $z = L$ (sobre las cuales la componente tangencial de $\mathbf{E}$ debe anularse), halle los valores discretos permitidos para $k$ y calcule el vector densidad superficial de corriente libre $\mathbf{K}_f(t)$ inducido sobre la placa en $z = 0$.
\end{enumerate}

\newpage

\section*{Solución Problema 1:}

\subsection*{a) Vectores $\mathbf{D}$, $\mathbf{E}$ y $\mathbf{P}$ por regiones}
\begin{figure}[H]
\centering
\begin{tikzpicture}[repaso,scale=1.04]
  \node at (0,2.16) {(a) Superficie gaussiana};
  \fill[black!13] (0,0) circle (1.44);
  \fill[black!32] (0,0) circle (.71);
  \draw[borde] (0,0) circle (.71) (0,0) circle (1.44);
  \draw[guia] (0,0) circle (1.08);
  \fill (0,0) circle (1.25pt) node[below right=-1pt] {$O$};
  \draw[eje] (0,0)--(181:.71) node[pos=.67,below=1pt] {$a$};
  \draw[eje] (0,0)--(222:1.44)
    node[pos=.83,above left=-1pt] {$b$};
  \draw[eje] (0,0)--(49:1.08)
    node[pos=.75,above left=-1pt] {$r$};
  \draw[vector] (7:1.08)--(7:1.79) node[right=1pt] {$\mathrm d\mathbf a$};
  \node[above=1pt] at (-.51,1.49) {$S_r$};
  \draw[llamada] (-.51,1.49)--(114:1.08);
  \node[anchor=east] at (-1.72,.46) {$+Q$ en $r=a$};
  \draw[llamada] (-1.63,.46)--(135:.71);
  \node[anchor=east] at (-1.70,-.38) {$\varepsilon$};
  \draw[llamada] (-1.62,-.38)--(197:1.26);
  \node[anchor=east] at (-1.70,-1.03) {$\varepsilon_0$};

  \begin{scope}[xshift=6.3cm]
    \node at (0,2.16) {(b) Polarización y cargas ligadas};
    \fill[black!13] (0,0) circle (1.44);
    \fill[black!32] (0,0) circle (.71);
    \draw[borde] (0,0) circle (.71) (0,0) circle (1.44);
    \foreach \ang in {27,72,...,342}{
      \node[inner sep=0pt] at (\ang:.84) {$-$};
      \node[inner sep=0pt] at (\ang:1.29) {$+$};
    }
    \foreach \ang in {6,51,96,186,231,276,321}{
      \draw[vector] (\ang:.93)--(\ang:1.22);
    }
    \draw[vector] (141:.93)--(141:1.22);
    \node[anchor=east] at (-1.76,.86) {$\mathbf P$};
    \draw[llamada] (-1.69,.81)--(141:1.10);
    \draw[vector] (108:.71)--(108:.22)
      node[below left=-1pt] {$\hat{\mathbf n}_a$};
    \draw[vector] (49:1.44)--(49:1.94)
      node[above right=-2pt] {$\hat{\mathbf n}_b$};
    \node at (0,-1.85) {$\hat{\mathbf n}_a=-\hat{\mathbf r},
        \qquad\hat{\mathbf n}_b=+\hat{\mathbf r}$};
  \end{scope}
\end{tikzpicture}
\caption{Corte central de la esfera. El gris oscuro identifica el
conductor y el gris claro, el dieléctrico. En (a), la esfera de Gauss
$S_r$ se dibuja en la corteza y su área vectorial apunta hacia afuera.
En (b), para $\varepsilon>\varepsilon_0$, los signos indican la carga
ligada de cada interfaz. Ambas normales salen del dieléctrico.}
\label{fig:e2-esfera}
\end{figure}

Por simetría esférica, los campos son radiales y dependen solo de $r$.
En la esfera gaussiana $S_r$ de la figura~\ref{fig:e2-esfera},
$\mathbf D=D_r(r)\hat{\mathbf r}$ y la normal exterior es
$\hat{\mathbf r}$. El elemento de área vectorial es
\begin{equation}
    d\mathbf a=\hat{\mathbf r}\,r^2\sin\theta\,d\theta\,d\phi.
    \label{eq:e2-area-esferica}
\end{equation}
Sustituimos \eqref{eq:e2-area-esferica} en el flujo de $\mathbf D$.
Sobre $S_r$, el radio es fijo, así que $D_r(r)$ sale de la integral.
\begin{align*}
    Q_{f,\mathrm{enc}}
    &=\oint_{S_r}\mathbf D\cdot d\mathbf a\\
    &=D_r(r)\int_0^{2\pi}\!\int_0^\pi
      \underbrace{(\hat{\mathbf r}\cdot\hat{\mathbf r})}_{1}
        r^2\sin\theta\,d\theta\,d\phi\\
    &=D_r(r)\,r^2(2\pi)\bigl[-\cos\theta\bigr]_0^\pi\\
    &=4\pi r^2D_r(r).
\end{align*}
En cada región usamos la carga libre encerrada por esa esfera.
\begin{itemize}
    \item \textbf{Interior del conductor ($r < a$):} En equilibrio electrostático el campo interno de un conductor es nulo y toda la carga libre $+Q$ reside en su superficie $r = a$. Por tanto:
    \begin{equation}
        \mathbf{E} = \mathbf{0}, \qquad \mathbf{D} = \mathbf{0}, \qquad \mathbf{P} = \mathbf{0}
    \end{equation}
    \item \textbf{Corteza dieléctrica ($a < r < b$):} La carga libre encerrada es $+Q$, de modo que $D(4\pi r^2) = Q$. Usando $\mathbf{D} = \varepsilon \mathbf{E}$ y $\mathbf{P} = \mathbf{D} - \varepsilon_0 \mathbf{E}$:
    \begin{equation}
        \mathbf{D}(r) = \frac{Q}{4\pi r^2} \hat{\mathbf{r}}, \qquad \mathbf{E}(r) = \frac{Q}{4\pi \varepsilon r^2} \hat{\mathbf{r}}, \qquad \mathbf{P}(r) = \frac{Q}{4\pi r^2} \left( 1 - \frac{\varepsilon_0}{\varepsilon} \right) \hat{\mathbf{r}}
    \end{equation}
    \item \textbf{Vacío exterior ($r > b$):} La carga libre encerrada sigue siendo $+Q$ y $\varepsilon = \varepsilon_0$:
    \begin{equation}
        \mathbf{D}(r) = \frac{Q}{4\pi r^2} \hat{\mathbf{r}}, \qquad \mathbf{E}(r) = \frac{Q}{4\pi \varepsilon_0 r^2} \hat{\mathbf{r}}, \qquad \mathbf{P}(r) = \mathbf{0}
    \end{equation}
\end{itemize}

\subsection*{b) Potencial del conductor, capacitancia y casos límite}
Para llegar desde el infinito hasta la superficie del conductor
seguimos un camino radial. Usamos $d\mathbf l=\hat{\mathbf r}\,dr$
con límites decrecientes. El campo cambia de expresión en $r=b$,
por lo que separamos allí el recorrido.
\begin{align*}
    V(a)-V(\infty)
    &=-\int_\infty^a\mathbf E\cdot d\mathbf l\\
    &=-\int_\infty^b \frac{Q}{4\pi\varepsilon_0 r^2}\,dr
      -\int_b^a \frac{Q}{4\pi\varepsilon r^2}\,dr\\
    &=-\frac{Q}{4\pi\varepsilon_0}
      \left[-\frac1r\right]_\infty^b
      -\frac{Q}{4\pi\varepsilon}
      \left[-\frac1r\right]_b^a\\
    &=\frac{Q}{4\pi\varepsilon_0 b}
      +\frac{Q}{4\pi\varepsilon}\left(\frac1a-\frac1b\right).
\end{align*}
Con $V(\infty)=0$, reagrupamos,
\[
    V(a)=\frac{Q}{4\pi}
      \left[\frac{1}{\varepsilon a}
        +\frac1b\left(\frac1{\varepsilon_0}-\frac1\varepsilon\right)\right].
\]
La capacitancia del sistema aislado es $C = Q / V(a)$:
\begin{equation}
    C = \frac{4\pi}{\dfrac{1}{\varepsilon a} + \dfrac{1}{b}\left( \dfrac{1}{\varepsilon_0} - \dfrac{1}{\varepsilon} \right)}
\end{equation}
Verificamos ambos límites físicos:
\begin{itemize}
    \item Si $\varepsilon \to \varepsilon_0$ (ausencia de dieléctrico): el segundo término del denominador da $\frac{1}{b}(0) = 0$, recuperando $C = 4\pi \varepsilon_0 a$, la capacitancia clásica de una esfera en el vacío.
    \item Si $b \to \infty$ (dieléctrico infinito): el término proporcional a $1/b$ se anula, entregando $C = 4\pi \varepsilon a$, correspondiente a una esfera inmersa en un medio dieléctrico ilimitado.
\end{itemize}

\subsection*{c) Cargas ligadas y verificación de apantallamiento}
En el volumen $a<r<b$, escribimos $\mathbf P=P_r(r)\hat{\mathbf r}$.
En coordenadas esféricas, $r^2P_r(r)$ es proporcional al flujo
de $\mathbf P$ por una esfera concéntrica. Al sustituir,
\begin{align*}
    \rho_b(r)
    &=-\nabla\cdot\mathbf P\\
    &=-\frac1{r^2}\frac{d}{dr}\bigl[r^2P_r(r)\bigr]\\
    &=-\frac1{r^2}\frac{d}{dr}
       \left[\frac{Q}{4\pi}\left(1-\frac{\varepsilon_0}{\varepsilon}\right)\right]\\
    &=0.
\end{align*}
Para las cargas superficiales usamos $\sigma_b=\mathbf P\cdot\hat{\mathbf n}$.
La figura~\ref{fig:e2-esfera}(b) muestra las dos normales que salen
del dieléctrico. Evaluamos $\mathbf P$ desde el material en cada caso.
\begin{align*}
    \sigma_b(a)
    &=P_r(a^+)\hat{\mathbf r}\cdot(-\hat{\mathbf r})\\
    &=-\frac{Q}{4\pi a^2}\left(1-\frac{\varepsilon_0}{\varepsilon}\right),\\[4pt]
    \sigma_b(b)
    &=P_r(b^-)\hat{\mathbf r}\cdot\hat{\mathbf r}\\
    &=+\frac{Q}{4\pi b^2}\left(1-\frac{\varepsilon_0}{\varepsilon}\right).
\end{align*}
Cada densidad es uniforme en su esfera. Por eso sale de la integral
y se multiplica por el área de la interfaz.
\begin{align*}
    Q_b(a)
    &=\oint_{S_a}\sigma_b(a)\,da\\
    &=4\pi a^2\,\sigma_b(a)=-Q\left(1-\frac{\varepsilon_0}{\varepsilon}\right),\\[4pt]
    Q_b(b)
    &=\oint_{S_b}\sigma_b(b)\,da\\
    &=4\pi b^2\,\sigma_b(b)=+Q\left(1-\frac{\varepsilon_0}{\varepsilon}\right).
\end{align*}
La carga ligada neta del cascarón es $Q_{b,\text{tot}} = Q_b(a) + Q_b(b) + \int \rho_b\, dV = 0$, como exige la neutralidad global del dieléctrico. Además, para una esfera gaussiana de radio $r \in (a, b)$, la carga total encerrada (libre más ligada) es:
\begin{equation}
    Q_{\text{enc}} = Q + Q_b(a) = Q - Q\left( 1 - \frac{\varepsilon_0}{\varepsilon} \right) = Q \frac{\varepsilon_0}{\varepsilon}
\end{equation}
Aplicando la Ley de Gauss fundamental $\oint \mathbf{E}\cdot d\mathbf{a} = Q_{\text{enc}}/\varepsilon_0$, resulta $E(r) = \frac{Q}{4\pi \varepsilon r^2}$, explicando microscópicamente la reducción del campo eléctrico dentro del dieléctrico.

\vspace{0.4cm}
\hrule
\vspace{0.5cm}

\section*{Solución Problema 2:}

\subsection*{a) Densidad de corriente superficial}
En el manto, el vector que va desde el eje a un punto de la pared
es $R\hat{\mathbf s}$. Con $\boldsymbol\omega=\omega\hat{\mathbf z}$,
la velocidad tangencial resulta
\begin{align*}
    \mathbf v
    &=\boldsymbol\omega\times(R\hat{\mathbf s})\\
    &=\omega R(\hat{\mathbf z}\times\hat{\mathbf s})\\
    &=\omega R\,\hat{\boldsymbol\phi}.
\end{align*}
Multiplicando por la densidad superficial de carga, obtenemos
\[
    \mathbf K=\sigma\mathbf v=\sigma\omega R\,\hat{\boldsymbol\phi}.
\]
La figura~\ref{fig:e2-cilindro-ampere}(a) muestra ese sentido
para $\sigma>0$ y $\omega>0$.
\subsection*{b) Campo magnético y condición de salto}
\begin{figure}[H]
\centering
\begin{tikzpicture}[repaso]
  \node at (0,3.04) {(a) Vista desde $+z$};
  \draw[borde] (0,1.24) circle (1.04);
  \node at (0,1.24) {$\odot$};
  \node[below=2pt] at (0,1.18) {$+\hat{\mathbf z}$};
  \draw[eje] (0,1.24)--(1.04,1.24)
    node[pos=.57,above=2pt] {$R$};
  \draw[vector] (1.04,1.24)--(1.04,2.03)
    node[right=2pt] {$\mathbf K$};
  \draw[vector] (18:1.43) ++(0,1.24)
    arc[start angle=18,end angle=112,radius=1.43];
  \node[above left=-1pt] at (-.46,2.58) {$\boldsymbol{\omega}$};

  \foreach \dx/\izq/\caso in {2.50/1.49/1,6.72/.59/2}{
    \begin{scope}[xshift=\dx cm]
      \fill[black!10] (\izq,.37) rectangle (2.78,2.30);
      \draw[eje] (0,0)--(3.58,0) node[right=1pt] {$s$};
      \draw[eje] (0,0)--(0,2.57) node[above=-1pt] {$z$};
      \draw[borde] (1.04,-.09)--(1.04,2.55);
      \node[below=2pt] at (1.04,-.09) {$s=R$};
      \foreach \zz in {.54,.93,1.33,1.73,2.13}
        \node[corte] at (1.04,\zz) {$\otimes$};
      \draw[guia] (\izq,.37) rectangle (2.78,2.30);
      \pgfmathsetmacro{\medio}{(\izq+2.78)/2}
      \draw[trayecto] (\izq,1.03)--(\izq,1.69);
      \draw[trayecto] (2.78,1.69)--(2.78,1.03);
      \draw[trayecto] ({\medio-.33},2.30)--({\medio+.33},2.30);
      \draw[trayecto] ({\medio+.33},.37)--({\medio-.33},.37);
      \ifnum\caso=1
        \node[right=1pt] at (\izq,1.31) {$\mathcal C_1$};
      \else
        \node[left=1pt] at (\izq,1.31) {$\mathcal C_1$};
      \fi
      \node[below=1pt] at (\medio,2.30) {$\mathcal C_2$};
      \node[left=1pt] at (2.78,1.65) {$\mathcal C_3$};
      \node[above=1pt] at (\medio,.37) {$\mathcal C_4$};
      \draw[medida] (3.18,.37)--(3.18,2.30)
        node[midway,right=1pt] {$h$};
      \ifnum\caso=1
        \node at (1.70,3.04) {(b) Lazo exterior};
        \node[above=1pt] at (\izq,2.30) {$s_1$};
        \node[above=1pt] at (2.78,2.30) {$s_2$};
        \node at (2.11,-.74) {$I_{\mathrm{enc}}=0$};
      \else
        \node at (1.70,3.04) {(c) Lazo que cruza el manto};
        \node[above=1pt] at (\izq,2.30) {$s<R$};
        \node[above=1pt] at (2.78,2.30) {$s_{\mathrm{ext}}>R$};
        \node at (2.11,-.74) {$I_{\mathrm{enc}}=Kh$};
      \fi
    \end{scope}
  }
\end{tikzpicture}
\caption{Cilindro cargado en rotación para $\omega>0$ y $\sigma>0$.
En (a), $\mathbf K$ es tangente al manto. En los cortes $sz$ de (b)
y (c), la corriente azimutal entra en la página ($\otimes$).
El recorrido sube en $\mathcal C_1$ y baja en $\mathcal C_3$.
En $\mathcal C_2$ y $\mathcal C_4$, el recorrido radial es
perpendicular al campo axial. Solo el lazo (c) corta el manto.}
\label{fig:e2-cilindro-ampere}
\end{figure}

El sistema es equivalente a un solenoide ideal infinito con corriente azimutal uniforme. Por invarianza rotacional en $\phi$, traslacional en $z$ y simetría de reflexión en planos transversales al eje $z$, el campo magnético tiene únicamente componente axial dependiente de $s$: $\mathbf{B}(s) = B_z(s) \hat{\mathbf{z}}$.
Para aplicar la ley de Ampère, escogemos los lazos de altura $h$
de la figura~\ref{fig:e2-cilindro-ampere}. Sean $s_1$ y $s_2$
los radios de los lados izquierdo y derecho, respectivamente.
Con $z'=0$ abajo y $z'=h$ arriba, los límites ya incluyen el sentido
de cada flecha. Partimos de
\[
    \oint_{\mathcal C}\mathbf B\cdot d\mathbf l
    =\sum_{i=1}^4\int_{\mathcal C_i}\mathbf B\cdot d\mathbf l
    =\mu_0 I_{\mathrm{enc}}.
\]
Usamos $d\mathbf l=\hat{\mathbf z}\,dz'$ en $\mathcal C_1$ y
$\mathcal C_3$, y $d\mathbf l=\hat{\mathbf s}\,ds'$ en $\mathcal C_2$
y $\mathcal C_4$. Sustituimos $\mathbf B=B_z\hat{\mathbf z}$ en
cada tramo.
\begin{align*}
    \int_{\mathcal C_1}\mathbf B\cdot d\mathbf l
    &=\int_0^h B_z(s_1)
          \underbrace{(\hat{\mathbf z}\cdot\hat{\mathbf z})}_{1}\,dz'\\
    &=hB_z(s_1),\\[4pt]
    \int_{\mathcal C_2}\mathbf B\cdot d\mathbf l
    &=\int_{s_1}^{s_2} B_z(s')
          \underbrace{(\hat{\mathbf z}\cdot\hat{\mathbf s})}_{0}\,ds'\\
    &=0,\\[4pt]
    \int_{\mathcal C_3}\mathbf B\cdot d\mathbf l
    &=\int_h^0 B_z(s_2)
          \underbrace{(\hat{\mathbf z}\cdot\hat{\mathbf z})}_{1}\,dz'\\
    &=-hB_z(s_2),\\[4pt]
    \int_{\mathcal C_4}\mathbf B\cdot d\mathbf l
    &=\int_{s_2}^{s_1} B_z(s')
          \underbrace{(\hat{\mathbf z}\cdot\hat{\mathbf s})}_{0}\,ds'\\
    &=0.
\end{align*}
Al sumar,
\begin{equation}
    h[B_z(s_1)-B_z(s_2)]=\mu_0 I_{\mathrm{enc}}.
    \label{eq:e2-ampere-cilindro}
\end{equation}
\begin{itemize}
    \item \textbf{Lazo exterior.} Con $R<s_1<s_2$, la superficie no
    corta el manto e $I_{\mathrm{enc}}=0$. De
    \eqref{eq:e2-ampere-cilindro}, $B_z(s_1)=B_z(s_2)$ para cualquier
    par de radios exteriores. Como $\mathbf B\to\mathbf0$ cuando
    $s\to\infty$, esa constante vale cero.
    \item \textbf{Lazo que cruza el manto.} Ahora $s_1=s<R$ y
    $s_2=s_{\mathrm{ext}}>R$. La superficie corta el manto en un
    segmento de altura $h$. La normal del lazo apunta en
    $+\hat{\boldsymbol\phi}$. Proyectamos $\mathbf K$ sobre ella,
    \[
        I_{\mathrm{enc}}
        =\int_0^h\mathbf K\cdot\hat{\boldsymbol\phi}\,dz'
        =\sigma\omega R\,h.
    \]
    Como $B_z(s_{\mathrm{ext}})=0$, en \eqref{eq:e2-ampere-cilindro}
    solo queda el aporte de $\mathcal C_1$. Al cancelar $h$,
    \[
        \mathbf B_{\mathrm{int}}(s)
        =\mu_0\sigma\omega R\,\hat{\mathbf z},
        \qquad s<R.
    \]
\end{itemize}
Para verificar el salto en $s = R$, elegimos la normal apuntando del interior al exterior ($\hat{\mathbf{n}} = +\hat{\mathbf{s}}$):
\begin{equation}
    \mathbf{B}_{\text{ext}} - \mathbf{B}_{\text{int}} = \mathbf{0} - \mu_0 \sigma \omega R \hat{\mathbf{z}} = -\mu_0 \sigma \omega R \hat{\mathbf{z}}
\end{equation}
Evaluando el miembro derecho de la condición de frontera:
\begin{equation}
    \mu_0 (\mathbf{K} \times \hat{\mathbf{n}}) = \mu_0 (\sigma \omega R \hat{\boldsymbol{\phi}} \times \hat{\mathbf{s}}) = -\mu_0 \sigma \omega R \hat{\mathbf{z}}
\end{equation}
Ambos miembros coinciden idénticamente.

\subsection*{c) Potencial vectorial, gauge de Coulomb y momento dipolar}
\begin{figure}[H]
\centering
\begin{tikzpicture}[repaso]
  \foreach \dx/\ss/\condicion/\area in {0/1.02/{s<R}/{\pi s^2},
                                       5.05/1.63/{s>R}/{\pi R^2}}{
    \begin{scope}[xshift=\dx cm]
      \pgfmathsetmacro{\rf}{min(1.32,\ss)}
      \fill[black!15] (0,0) circle (\rf);
      \draw[guia] (0,0) circle (1.32);
      \draw[borde] (0,0) circle (\ss);
      \draw[vector] (48:\ss) arc[start angle=48,end angle=102,radius=\ss];
      \node[anchor=west] at (43:{\ss+.22}) {$d\mathbf l$};
      \draw[eje] (0,0)--(0:\ss) node[pos=.55,above=2pt] {$s$};
      \draw[eje] (0,0)--(154:1.32) node[pos=.56,above=1pt] {$R$};
      \node at (0,-.45) {$\odot\;\mathbf B$};
      \node[anchor=west] at (318:{\ss+.19}) {$\mathcal C_s$};
      \node at (0,2.04) {$\condicion$};
      \node at (0,-1.99) {$A_{\mathrm{flujo}}=\area$};
    \end{scope}
  }
\end{tikzpicture}
\caption{Discos de Stokes vistos desde $+z$. El borde $\mathcal C_s$
se recorre en el sentido $+\hat{\boldsymbol\phi}$ y la normal
apunta hacia $+z$. El sombreado es la parte del disco que atraviesa
el campo magnético. Para todo contorno con $s>R$, el área que
aporta flujo sigue siendo $\pi R^2$.}
\label{fig:e2-cilindro-stokes}
\end{figure}

Por simetría, proponemos $\mathbf A=A_\phi(s)\hat{\boldsymbol\phi}$.
Aplicamos el teorema de Stokes al disco $S_s$ de la
figura~\ref{fig:e2-cilindro-stokes}, con normal $+\hat{\mathbf z}$ y
borde $\mathcal C_s$ orientado en $+\hat{\boldsymbol\phi}$.
\begin{align*}
    \oint_{\mathcal C_s}\mathbf A\cdot d\mathbf l
    &=\int_{S_s}\mathbf B\cdot d\mathbf a,\\
    \oint_{\mathcal C_s}\mathbf A\cdot d\mathbf l
    &=A_\phi(s)\int_0^{2\pi}
        (\hat{\boldsymbol\phi}\cdot\hat{\boldsymbol\phi})\,s\,d\phi\\
    &=2\pi s\,A_\phi(s).
\end{align*}
En el flujo usamos $d\mathbf a=\hat{\mathbf z}\,s'\,ds'\,d\phi$.
Como el campo se anula para $s'>R$, solo aporta el disco sombreado.
\begin{align*}
    \int_{S_s}\mathbf B\cdot d\mathbf a
    &=(\mu_0\sigma\omega R)
      \int_0^{2\pi}\!\int_0^{\min(s,R)}
       (\hat{\mathbf z}\cdot\hat{\mathbf z})\,s'\,ds'\,d\phi\\
    &=(\mu_0\sigma\omega R)\,\pi\min(s,R)^2.
\end{align*}
Igualamos la circulación y el flujo, y dividimos por $2\pi s$.
\begin{align*}
    \mathbf A_{\mathrm{int}}(s)
    &=\frac12\mu_0\sigma\omega R s\,\hat{\boldsymbol\phi},
      &&0<s<R,\\
    \mathbf A_{\mathrm{ext}}(s)
    &=\frac{\mu_0\sigma\omega R^3}{2s}\,\hat{\boldsymbol\phi},
      &&s>R.
\end{align*}
Ambas expresiones coinciden en $s=R$. La solución interior tiende
a $\mathbf0$ en el eje, como exige $\mathbf A(0)=\mathbf0$.
En coordenadas cilíndricas, $\nabla \cdot \mathbf{A} = \frac{1}{s}\frac{\partial A_\phi}{\partial \phi} = 0$, ya que $A_\phi$ solo depende de $s$; por ende, pertenece al gauge de Coulomb.

Finalmente, un anillo diferencial de longitud axial $dz$ transporta corriente $dI = K\, dz = \sigma \omega R\, dz$ alrededor del área $\pi R^2 \hat{\mathbf{z}}$. Su momento dipolar es $d\mathbf{m} = dI (\pi R^2 \hat{\mathbf{z}})$, de donde el momento dipolar magnético por unidad de longitud resulta:
\begin{equation}
    \frac{d\mathbf{m}}{dz} = \pi \sigma \omega R^3 \hat{\mathbf{z}}
\end{equation}

\vspace{0.4cm}
\hrule
\vspace{0.5cm}

\section*{Solución Problema 3:}

\subsection*{a) Campos $B$, $H_m$ y $H_0$ en el circuito magnético}
\begin{figure}[H]
\centering
\begin{tikzpicture}[repaso]
  \node at (0,2.58) {(a) Vista superior del toroide};
  \path[fill=black!14]
    (11:1.91) arc[start angle=11,end angle=349,radius=1.91]
    -- (349:1.13) arc[start angle=349,end angle=11,radius=1.13] -- cycle;
  \draw[borde] (11:1.91) arc[start angle=11,end angle=349,radius=1.91]
    -- (349:1.13) arc[start angle=349,end angle=11,radius=1.13] -- cycle;
  \draw[guia] (0,0) circle (1.52);
  \draw[trayecto] (45:1.52) arc[start angle=45,end angle=78,radius=1.52];
  \node at (69:2.07) {$\mathcal C_m$};
  \draw[vector] (349:1.52) arc[start angle=349,end angle=371,radius=1.52];
  \node[anchor=west] at (1.64,0) {$\mathcal C_0$};
  \draw[eje] (0,0)--(59:1.52) node[pos=.57,right=2pt] {$R$};
  \foreach \ang in {127,151,175}{
    \draw[llamada] (\ang:1.04)--(\ang:2.00);
    \node[corte] at (\ang:1.20) {$\odot$};
    \node[corte] at (\ang:1.84) {$\otimes$};
  }
  \foreach \ang in {230,254}{
    \draw[llamada] (\ang:1.03)--(\ang:2.00);
    \node[corte] at (\ang:1.20) {$\circ$};
    \node[corte] at (\ang:1.84) {$\circ$};
  }
  \node[anchor=east] at (-2.19,.86) {Primario $N$};
  \draw[llamada] (-2.09,.86)--(151:1.91);
  \node[anchor=east] at (-1.65,-1.76) {Secundario $N_2$};
  \draw[llamada] (-1.56,-1.72)--(230:1.91);

  \begin{scope}[xshift=4.36cm]
    \node at (2.12,2.58) {(b) Detalle del entrehierro};
    \fill[black!14] (0,-.95) rectangle (1.43,1.48)
      (2.79,-.95) rectangle (4.22,1.48);
    \draw[borde] (0,-.95) rectangle (1.43,1.48)
      (2.79,-.95) rectangle (4.22,1.48);
    \node at (.71,1.13) {$\mu$};
    \node at (2.11,1.37) {$\mu_0$};
    \node at (3.50,1.13) {$\mu$};
    \draw[llamada] (1.43,1.54)--(1.43,1.96)
      (2.79,1.54)--(2.79,1.96);
    \draw[medida] (1.43,1.85)--(2.79,1.85)
      node[midway,above=2pt] {$d$};
    \draw[vector] (.28,.63)--(3.94,.63)
      node[midway,above=2pt] {$\mathbf B$};
    \draw[vector] (.48,-.04)--(.94,-.04)
      node[midway,above=1pt] {$\mathbf H_m$};
    \draw[vector] (1.60,-.04)--(2.62,-.04)
      node[midway,above=1pt] {$\mathbf H_0$};
    \draw[vector] (3.28,-.04)--(3.74,-.04)
      node[midway,above=1pt] {$\mathbf H_m$};
    \draw[vector] (.36,-.68)--(1.08,-.68)
      node[midway,above=1pt] {$\mathbf M$};
    \node[corte] at (.75,-.95) {$\otimes$};
    \node at (1.00,-1.32) {$\mathbf K_b$ en el manto};
  \end{scope}
\end{tikzpicture}
\caption{Toroide con entrehierro. En (a), el camino amperiano va por
la circunferencia media y se separa en $\mathcal C_m$ (material) y
$\mathcal C_0$ (vacío). El primario atraviesa el disco amperiano por
el lado del hueco, en $+z$ ($\odot$). En (b), $\mathbf B$ cruza
normalmente las caras de la ranura. Para $\mu>\mu_0$, el campo
$\mathbf H$ es mayor en el vacío y $\mathbf M$ existe en el material.
En el borde inferior del dibujo se indica un sentido de $\mathbf K_b$.}
\label{fig:e2-toroide}
\end{figure}

Por simetría toroidal y despreciando efectos de borde en la ranura estrecha ($d \ll R$), las líneas de $\mathbf{B}$ siguen circunferencias azimutales concéntricas e inciden perpendicularmente sobre las dos caras planas que delimitan el entrehierro. Como la componente normal de $\mathbf{B}$ es siempre continua ($\Delta B^\perp = 0$), el módulo $B$ en el material es igual al módulo en la ranura de vacío: $B_m = B_0 \equiv B$.

Usando las relaciones constitutivas en cada tramo, los módulos del campo auxiliar son:
\begin{equation}
    H_m = \frac{B}{\mu} \quad (\text{en la longitud } 2\pi R - d), \qquad H_0 = \frac{B}{\mu_0} \quad (\text{en la longitud } d)
\end{equation}
Aplicamos la ley de Ampère sobre la circunferencia media de la
figura~\ref{fig:e2-toroide}(a). La superficie limitada por la
circunferencia es atravesada por las $N$ vueltas del primario
en el mismo sentido. Por eso $I_{f,\mathrm{enc}}=NI$.
En los dos tramos, $d\mathbf l=\hat{\boldsymbol\phi}\,dl$.
\begin{align*}
    NI
    &=\oint_{\mathcal C}\mathbf H\cdot d\mathbf l\\
    &=\int_{\mathcal C_m}\mathbf H_m\cdot d\mathbf l
       +\int_{\mathcal C_0}\mathbf H_0\cdot d\mathbf l\\
    &=H_m\int_{\mathcal C_m}
          (\hat{\boldsymbol\phi}\cdot\hat{\boldsymbol\phi})\,dl
      +H_0\int_{\mathcal C_0}
          (\hat{\boldsymbol\phi}\cdot\hat{\boldsymbol\phi})\,dl\\
    &=H_m(2\pi R-d)+H_0 d\\
    &=B\left(\frac{2\pi R-d}{\mu}+\frac{d}{\mu_0}\right).
\end{align*}
En el último paso usamos que el campo cruza perpendicularmente
las caras del entrehierro, como se ve en la
figura~\ref{fig:e2-toroide}(b). La continuidad de su componente
normal permite usar el mismo $B$ para calcular $H_m$ y $H_0$.
Despejando $B$, $H_m$ y $H_0$:
\begin{align}
    B &= \frac{\mu_0 \mu N I}{\mu_0(2\pi R - d) + \mu d} \\
    H_m &= \frac{\mu_0 N I}{\mu_0(2\pi R - d) + \mu d}, \qquad H_0 = \frac{\mu N I}{\mu_0(2\pi R - d) + \mu d} = (1 + \chi_m) H_m
\end{align}

\subsection*{b) Corriente superficial ligada y discontinuidad de $H$ en el entrehierro}
Dentro del núcleo, $\mathbf M=\chi_m\mathbf H_m$. En el manto
lateral, $\hat{\mathbf n}$ es perpendicular a $\mathbf M$.
En las caras del entrehierro, ambas direcciones son colineales.
Por tanto, la magnitud de la corriente ligada en el manto es
\begin{align*}
    K_b
    &=\bigl|\mathbf M\times\hat{\mathbf n}\bigr|\\
    &=|\chi_m|\,H_m\sin\left(\frac{\pi}{2}\right)\\
    &=\frac{|\chi_m|\mu_0 NI}{\mu_0(2\pi R-d)+\mu d}.
\end{align*}
En las caras del corte ese producto cruz se anula. La siguiente
interpretación corresponde al caso $\mu>\mu_0$ de la figura.
\textit{Explicación física del salto entre $H_m$ y $H_0$:} La ausencia de corrientes libres ($\mathbf{K}_f = \mathbf{0}$) solo garantiza la continuidad de la componente \textbf{tangencial} $\mathbf{H}^\parallel$. En las caras planas del entrehierro, el campo incide de manera \textbf{normal} a la interfaz. Dado que $\nabla \cdot \mathbf{H} = -\nabla \cdot \mathbf{M}$, al pasar abruptamente del material magnetizado ($M^\perp = M$) al vacío ($M^\perp = 0$), aplica la condición normal $\Delta H^\perp = -\Delta M^\perp$, esto es, $H_0 - H_m = -(0 - M) = M$. Esta acumulación efectiva de ``polos magnéticos'' en las caras del corte refuerza $H_0$ en el vacío y genera un campo desmagnetizante que reduce $H_m$ dentro del núcleo.

\subsection*{c) Inductancia mutua y fem inducida en el secundario}
Tomamos la sección transversal $S_A$ del núcleo con la normal en
el sentido de $\mathbf B$. En la aproximación del núcleo delgado,
$B$ es uniforme en esa sección.
\begin{align*}
    \Phi_1
    &=\int_{S_A}\mathbf B\cdot d\mathbf a\\
    &=B\int_{S_A}
       (\hat{\boldsymbol\phi}\cdot\hat{\boldsymbol\phi})\,da\\
    &=BA.
\end{align*}
Cada una de las $N_2$ vueltas de la bobina secundaria enlaza el
mismo flujo. Por tanto, $\Phi_2=N_2\Phi_1=\mathcal M I$.
\begin{align*}
    \mathcal M
    &=\frac{N_2BA}{I}\\
    &=\frac{\mu_0\mu NN_2 A}{\mu_0(2\pi R-d)+\mu d}.
\end{align*}
Si $I(t) = I_0 \cos(\omega t)$, por la Ley de Faraday la fuerza electromotriz inducida en el secundario es:
\begin{equation}
    \mathcal{E}_2(t) = -\mathcal{M} \frac{dI}{dt} = \frac{\mu_0 \mu N N_2 A I_0 \omega}{\mu_0(2\pi R - d) + \mu d} \sin(\omega t)
\end{equation}

\vspace{0.4cm}
\hrule
\vspace{0.5cm}

\section*{Solución Problema 4:}

\subsection*{a) Descarga radial y tiempo de relajación}
\begin{figure}[H]
\centering
\begin{tikzpicture}[repaso]
  \node at (0,2.22) {(a) Corte transversal};
  \fill[black!13] (0,0) circle (1.55);
  \fill[white] (0,0) circle (.60);
  \draw[borde] (0,0) circle (.60) (0,0) circle (1.55);
  \draw[guia] (0,0) circle (1.04);
  \fill (0,0) circle (1.15pt);
  \draw[eje] (0,0)--(142:.60) node[pos=.73,above=1pt] {$a$};
  \draw[eje] (0,0)--(222:1.55) node[pos=.78,left=1pt] {$b$};
  \draw[eje] (0,0)--(7:1.04) node[pos=.53,above=2pt] {$s$};
  \foreach \ang in {45,100,170,270,325}
    \draw[vector] (\ang:.72)--(\ang:1.37);
  \node[anchor=west] at (1.69,.75) {$S_s$};
  \draw[llamada] (1.64,.70)--(30:1.04);
  \node[anchor=east] at (-1.69,1.06) {$+\lambda(t)$};
  \draw[llamada] (-1.58,1.00)--(118:.60);
  \node[anchor=west] at (1.69,-.98) {$-\lambda(t)$};
  \draw[llamada] (1.59,-.94)--(327:1.55);
  \node at (0,-1.94) {$\mathbf D,\ \mathbf E,\ \mathbf J_f$
    apuntan hacia afuera};

  \begin{scope}[xshift=6.61cm]
    \node at (0,2.22) {(b) Elemento del manto de $S_s$};
    \path[fill=black!13,draw=black!60,line width=.5pt]
      (-.60,-.91)--(.20,-.50)--(.20,1.26)--(-.60,.85)--cycle;
    \draw[guia] (-.20,-.68)--(-.20,.92);
    \fill (-.20,.05) circle (1.2pt);
    \draw[vector] (-.20,.05)--(1.29,.05)
      node[right=2pt] {$\mathbf J_f$};
    \draw[vector] (-.20,.05)--(-1.69,.05)
      node[left=2pt] {$\mathbf J_d$};
    \draw[vector] (-.20,.72)--(1.13,.72)
      node[right=2pt] {$d\mathbf a\parallel+\hat{\mathbf s}$};
    \node at (0,-1.42) {$\mathbf J_f\cdot d\mathbf a
                          +\mathbf J_d\cdot d\mathbf a=0$};
  \end{scope}
\end{tikzpicture}
\caption{Descarga radial para $\lambda(t)>0$. En (a), la línea
discontinua representa el corte del cilindro gaussiano de radio
$s$ y longitud $\ell$. En (b), las dos corrientes tienen igual
magnitud y sentidos opuestos en el mismo elemento del manto.}
\label{fig:e2-descarga}
\end{figure}

Sea $\lambda(t)$ la carga libre por unidad de longitud del
conductor interno. Elegimos la superficie gaussiana $S_s$
de la figura~\ref{fig:e2-descarga}, con radio $a<s<b$ y longitud $\ell$.
Por simetría, $\mathbf D=D_s(s,t)\hat{\mathbf s}$. Las tapas tienen
normal $\pm\hat{\mathbf z}$. En el manto,
$d\mathbf a=\hat{\mathbf s}\,s\,d\phi\,dz$.
\begin{align*}
    \lambda(t)\ell
    &=\oint_{S_s}\mathbf D\cdot d\mathbf a\\
    &=\int_{\mathrm{manto}}D_s(s,t)
        \underbrace{(\hat{\mathbf s}\cdot\hat{\mathbf s})}_{1}\,da
     +\int_{\mathrm{tapas}}\mathbf D\cdot d\mathbf a\\
    &=D_s(s,t)\int_0^\ell\!\int_0^{2\pi}s\,d\phi\,dz+0\\
    &=D_s(s,t)(2\pi s\ell).
\end{align*}
El segundo aporte se anula porque
$\hat{\mathbf s}\cdot(\pm\hat{\mathbf z})=0$.
Dividimos por el área del manto. Después usamos
$\mathbf E=\mathbf D/\varepsilon$ y $\mathbf J_f=\sigma\mathbf E$.
\begin{align*}
    \mathbf D(s,t)&=\frac{\lambda(t)}{2\pi s}\hat{\mathbf s},\\
    \mathbf E(s,t)&=\frac{\lambda(t)}{2\pi\varepsilon s}\hat{\mathbf s},\\
    \mathbf J_f(s,t)&=\frac{\sigma\lambda(t)}{2\pi\varepsilon s}
                       \hat{\mathbf s}.
\end{align*}
Para la corriente libre que sale del cilindro usamos el mismo
elemento de área vectorial del manto,
\begin{align*}
    I_{f,\mathrm{rad}}
    &=\int_{\mathrm{manto}}\mathbf J_f\cdot d\mathbf a\\
    &=J_f(s,t)(2\pi s\ell)\\
    &=\frac{\sigma}{\varepsilon}\lambda(t)\ell.
\end{align*}
La carga libre encerrada disminuye al ritmo al que sale la corriente,
\begin{align*}
    \frac{d[\lambda(t)\ell]}{dt}&=-I_{f,\mathrm{rad}},\\
    -\frac{d\lambda}{dt}&=\frac{\sigma}{\varepsilon}\lambda(t).
\end{align*}
Con $\lambda(0)=\lambda_0$ resulta
$\lambda(t)=\lambda_0 e^{-t/\tau}$, con $\tau=\varepsilon/\sigma$.
Sustituyendo $\lambda(t)$, los campos explícitos para $t > 0$ son:
\begin{equation}
    \mathbf{D}(s,t) = \frac{\lambda_0 e^{-t/\tau}}{2\pi s} \hat{\mathbf{s}}, \qquad \mathbf{E}(s,t) = \frac{\lambda_0 e^{-t/\tau}}{2\pi \varepsilon s} \hat{\mathbf{s}}, \qquad \mathbf{J}_f(s,t) = \frac{\lambda_0 e^{-t/\tau}}{2\pi \tau s} \hat{\mathbf{s}}
\end{equation}

\subsection*{b) Cancelación exacta con la corriente de desplazamiento}
En el medio, la corriente de desplazamiento se obtiene derivando
el desplazamiento eléctrico a radio fijo.
\begin{align*}
    \mathbf J_d(s,t)
    &=\frac{\partial\mathbf D}{\partial t}\\
    &=\frac{\hat{\mathbf s}}{2\pi s}
       \frac{d}{dt}\bigl[\lambda_0e^{-t/\tau}\bigr]\\
    &=-\frac{\lambda_0e^{-t/\tau}}{2\pi\tau s}\hat{\mathbf s}\\
    &=-\mathbf J_f(s,t).
\end{align*}
Sobre el elemento del manto de la figura~\ref{fig:e2-descarga}(b),
proyectamos ambas densidades sobre la misma normal exterior,
\[
    \mathbf J_f\cdot d\mathbf a=J_f\,da,
    \qquad
    \mathbf J_d\cdot d\mathbf a=-J_f\,da.
\]
La densidad de corriente libre de conducción y la densidad de corriente de desplazamiento macroscópica son iguales en magnitud y exactamente opuestas en cada punto del medio ($\mathbf{J}_f + \mathbf{J}_d = \mathbf{0}$).
Para cualquier superficie delimitada por $\mathcal C$, la ley
de Ampère--Maxwell da
\begin{align*}
    \oint_{\mathcal C}\mathbf H\cdot d\mathbf l
    &= I_{f,\mathrm{enc}}+\frac{d\Phi_D}{dt}\\
    &=\int_S(\mathbf J_f+\mathbf J_d)\cdot d\mathbf a\\
    &=0.
\end{align*}
Para el campo inducido, la simetría cilíndrica y la reflexión
en cualquier plano que contiene el eje permiten solo una componente
azimutal. Sobre una circunferencia coaxial de radio $s$,
\[
    \oint_{\mathcal C_s}\mathbf H\cdot d\mathbf l
    =H_\phi(s,t)\int_0^{2\pi}s\,d\phi=2\pi s\,H_\phi(s,t).
\]
El resultado anterior exige que esta circulación sea cero a todo
radio. Por tanto, $\mathbf H=\mathbf0$ y $\mathbf B=\mathbf0$
en todo el espacio durante la descarga.


\vspace{0.4cm}
\hrule
\vspace{0.5cm}

\section*{Solución Problema 5: Síntesis de las Ecuaciones de Maxwell en la Materia}

\subsection*{a) Divergencia magnética, campo $\mathbf{H}$ y magnetización $\mathbf{M}$}
La única componente de $\mathbf B$ es $B_y(z,t)$ y su amplitud
cambia al avanzar en $z$. La divergencia deriva cada componente
respecto de su propia coordenada.
\begin{align*}
    \nabla\cdot\mathbf B
    &=\frac{\partial B_x}{\partial x}
      +\frac{\partial B_y}{\partial y}
      +\frac{\partial B_z}{\partial z}\\
    &=0+\frac{\partial}{\partial y}
       [B_0\cos(kz)\sin(\omega t)]+0\\
    &=0.
\end{align*}
Como $\mu = \mu_r \mu_0$ y $\mathbf{B} = \mu_0(\mathbf{H} + \mathbf{M}) = \mu \mathbf{H}$:
\begin{equation}
    \mathbf{H}(z,t) = \frac{B_0}{\mu_r \mu_0} \cos(kz) \sin(\omega t) \hat{\mathbf{y}}, \qquad \mathbf{M}(z,t) = \frac{(\mu_r - 1) B_0}{\mu_r \mu_0} \cos(kz) \sin(\omega t) \hat{\mathbf{y}}
\end{equation}

\subsection*{b) Campos $\mathbf{D}$, $\mathbf{E}$, polarización $\mathbf{P}$ y Ley de Gauss}
Como $H_y$ depende de $z$, el único término del rotacional es el
que contiene $\hat{\mathbf z}\times\hat{\mathbf y}=-\hat{\mathbf x}$.
\begin{align*}
    \nabla\times\mathbf H
    &=(\hat{\mathbf z}\times\hat{\mathbf y})
        \frac{\partial H_y}{\partial z}\\
    &=-\hat{\mathbf x}\frac{\partial}{\partial z}
        \left[\frac{B_0}{\mu_r\mu_0}\cos(kz)\sin(\omega t)\right]\\
    &=\frac{kB_0}{\mu_r\mu_0}\sin(kz)\sin(\omega t)\hat{\mathbf x}.
\end{align*}
Por Ampère--Maxwell sin corriente libre,
$\partial_t\mathbf D=\nabla\times\mathbf H$. Integramos a $z$
fijo y tomamos nulo el término electrostático.
\begin{align*}
    \mathbf D(z,t)
    &=\frac{kB_0}{\mu_r\mu_0}\sin(kz)
          \left(\int\sin(\omega t)\,dt\right)\hat{\mathbf x}\\
    &=-\frac{kB_0}{\omega\mu_r\mu_0}
        \sin(kz)\cos(\omega t)\hat{\mathbf x}.
\end{align*}
Dividiendo por $\varepsilon = \varepsilon_r \varepsilon_0$ y recordando que $\varepsilon_0 \mu_0 = 1/c^2$, hallamos $\mathbf{E}$ y $\mathbf{P} = \mathbf{D} - \varepsilon_0 \mathbf{E}$:
\begin{align}
    \mathbf{E}(z,t) &= -\frac{c^2 k B_0}{\omega \varepsilon_r \mu_r} \sin(kz) \cos(\omega t) \hat{\mathbf{x}} \\
    \mathbf{P}(z,t) &= -\frac{(\varepsilon_r - 1) k B_0}{\varepsilon_r \omega \mu_r \mu_0} \sin(kz) \cos(\omega t) \hat{\mathbf{x}}
\end{align}
Como $D_x$ depende únicamente de $z$ y $t$, se verifica $\nabla \cdot \mathbf{D} = \frac{\partial D_x}{\partial x} = 0$, cumpliendo la Ley de Gauss en ausencia de carga libre.

\subsection*{c) Relación de dispersión, modos confinados y corriente superficial en $z = 0$}
Partimos de la ley de Faraday,
$\nabla\times\mathbf E=-\partial\mathbf B/\partial t$.
Como $\mathbf E=E_x(z,t)\hat{\mathbf x}$, la derivada en $z$
aporta en el sentido $\hat{\mathbf z}\times\hat{\mathbf x}=\hat{\mathbf y}$.
\begin{align*}
    \nabla\times\mathbf E
    &=\hat{\mathbf y}\frac{\partial E_x}{\partial z}\\
    &=-\frac{c^2 k^2 B_0}{\omega\varepsilon_r\mu_r}
         \cos(kz)\cos(\omega t)\hat{\mathbf y},\\[4pt]
    -\frac{\partial\mathbf B}{\partial t}
    &=-B_0\cos(kz)\frac{d}{dt}\sin(\omega t)\hat{\mathbf y}\\
    &=-\omega B_0\cos(kz)\cos(\omega t)\hat{\mathbf y}.
\end{align*}
Igualamos los coeficientes de $\cos(kz)\cos(\omega t)\hat{\mathbf y}$
para que la ley se cumpla en todo $z$ y $t$.
\begin{align*}
    \frac{c^2 k^2}{\omega\varepsilon_r\mu_r}&=\omega,\\
    \omega^2&=\frac{c^2 k^2}{\varepsilon_r\mu_r},\\
    \omega&=\frac{ck}{\sqrt{\varepsilon_r\mu_r}}.
\end{align*}
\begin{figure}[H]
\centering
\begin{tikzpicture}[repaso]
  \node at (2.40,2.30) {(a) Modo $m=1$ entre las placas};
  \fill[black!14] (-.13,-1.21) rectangle (0,1.24)
    (4.79,-1.21) rectangle (4.92,1.24);
  \draw[borde] (0,-1.21)--(0,1.24) (4.79,-1.21)--(4.79,1.24);
  \draw[eje] (-.32,0)--(5.23,0) node[right=1pt] {$z$};
  \draw[line width=.88pt,domain=0:4.79,samples=55,variable=\u]
    plot (\u,{-1.04*sin(180*\u/4.79)});
  \draw[black!72,line width=.80pt,densely dashed,
        domain=0:4.79,samples=55,variable=\u]
    plot (\u,{1.04*cos(180*\u/4.79)});
  \fill (0,0) circle (1.6pt) (4.79,0) circle (1.6pt);
  \node at (0,-1.51) {$0$};
  \node at (4.79,-1.51) {$L$};
  \node[above right] at (.56,.93) {$H_y/H_*$};
  \node[below] at (2.36,-1.03) {$E_x/E_*$};
  \node at (2.40,1.73) {$\omega t=\pi/4$};

  \begin{scope}[xshift=7.1cm]
    \node at (.69,2.30) {(b) Placa en $z=0$};
    \fill[black!17] (-1.02,-1.25) rectangle (0,1.45);
    \draw[borde] (0,-1.25)--(0,1.45);
    \node at (1.23,1.38) {Medio};
    \draw[vector] (0,.15)--(1.18,.15)
      node[right=1pt] {$\hat{\mathbf n}=+\hat{\mathbf z}$};
    \draw[vector] (0,.15)--(0,-.86);
    \node[anchor=west] at (.15,-.65) {$\mathbf K_f\parallel-\hat{\mathbf x}$};
    \node[anchor=west] at (.27,.82) {$\odot\;\mathbf H(0^+,t)$};
    \node at (.54,-1.58) {$E_x(0,t)=0$};
  \end{scope}
\end{tikzpicture}
\caption{Onda estacionaria en el modo fundamental. En (a), cada
componente se divide por su propia amplitud máxima en el instante
indicado. El campo eléctrico tangencial tiene nodos en ambas
placas. En (b), $+x$ apunta hacia arriba y $+y$ sale de la página,
de modo que $\hat{\mathbf z}\times\mathbf H(0^+,t)$ apunta hacia
$-\hat{\mathbf x}$ en ese instante.}
\label{fig:e2-onda}
\end{figure}

Las placas de la figura~\ref{fig:e2-onda}(a) son paralelas a
$xy$. Como $\mathbf E$ apunta en $\hat{\mathbf x}$, todo el campo
eléctrico es tangencial a sus caras. Por tanto, exigimos
$E_x(0,t)=E_x(L,t)=0$.
\begin{itemize}
    \item En $z = 0$: $E_x(0,t) \propto \sin(0) = 0$ se satisface automáticamente.
    \item En $z = L$: $E_x(L,t) = 0$ con solución no trivial exige $\sin(kL) = 0$, lo que cuantiza los números de onda permitidos en:
    \begin{equation}
        k_m = \frac{m \pi}{L}, \qquad m = 1, 2, 3, \dots
    \end{equation}
\end{itemize}
En la placa de $z=0$ de la figura~\ref{fig:e2-onda}(b), la normal
del conductor hacia el medio es $+\hat{\mathbf z}$ y
$\mathbf H_1=\mathbf0$ dentro del conductor. Evaluamos el campo
del lado del medio usando $\cos(k\,0)=1$,
\[
    \mathbf H_2(0^+,t)=\frac{B_0}{\mu_r\mu_0}
                           \sin(\omega t)\hat{\mathbf y}.
\]
Partimos de la condición tangencial
$\mathbf H_2-\mathbf H_1=\mathbf K_f\times\hat{\mathbf n}$.
Multiplicamos vectorialmente por $\hat{\mathbf n}$ a la izquierda.
Como $\mathbf K_f$ es tangencial, $\hat{\mathbf n}\cdot\mathbf K_f=0$,
y el producto triple se reduce a
\[
    \hat{\mathbf n}\times(\mathbf K_f\times\hat{\mathbf n})
    =\mathbf K_f(\hat{\mathbf n}\cdot\hat{\mathbf n})
       -\hat{\mathbf n}(\hat{\mathbf n}\cdot\mathbf K_f)
    =\mathbf K_f.
\]
Sustituimos el campo y la normal,
\begin{align*}
    \mathbf K_f(t)
    &=\hat{\mathbf z}\times\mathbf H_2(0^+,t)\\
    &=\frac{B_0}{\mu_r\mu_0}\sin(\omega t)
       (\hat{\mathbf z}\times\hat{\mathbf y})\\
    &=-\frac{B_0}{\mu_r\mu_0}\sin(\omega t)\hat{\mathbf x}.
\end{align*}
(Se comprueba de inmediato que $\mathbf{K}_f \times \hat{\mathbf{z}} = (-\hat{\mathbf{x}} \times \hat{\mathbf{z}})\frac{B_0}{\mu_r \mu_0}\sin(\omega t) = +\frac{B_0}{\mu_r \mu_0}\sin(\omega t)\hat{\mathbf{y}} = \mathbf{H}_2(0^+,t)$).

\end{document}